\documentclass{article}

%
% This lecture is based on Shriram's notes for the first
%  Rice 2000 workshop.
%

\usepackage{latexsym}
\usepackage{slatex}
\defschememathescape{$}
\setspecialsymbol{->}{$\longrightarrow$}
\setspecialsymbol{...}{$\cdots$}
\setspecialsymbol{---}{$\Box$}
\setkeyword{local}
\parindent=0pt

\begin{document}

\begin{verbatim}
; Let's write quicksort... First, a data defn.

; A num-list is either
; - empty, or
; - (cons num num-list)

; Based on our knowledge of the algorithm, we know we need helper
; functions. Let's write them first.
\end{verbatim}

\begin{schemedisplay}
(define (nl-proc pivot nl)
  (cond
    [(empty? nl) ...]
    [(cons? nl)
     ... (first nl) ...
     ... (nl-proc pivot (rest nl)) ...]))
\end{schemedisplay}

\begin{verbatim}
; We need helper functions to split the list
;  for a pivot item

; nums-le : num num-list -> num-list
\end{verbatim}

\begin{schemedisplay}
(define (nums-le pivot nl)
  (cond
    [(empty? nl) ...]
    [(cons? nl)
     (cond
       [(<= pivot (first nl))
        ... (first nl) ...
        ... (nums-le pivot (rest nl)) ...]
       [else
        ... (first nl) ...
        ... (nums-le pivot (rest nl)) ...])]))

(define (nums-le pivot nl)
  (cond
    [(empty? nl) empty]
    [(cons? nl)
     (cond
       [(<= (first nl) pivot)
        (cons (first nl) (nums-le pivot (rest nl)))]
       [else
        (nums-le pivot (rest nl))])]))
\end{schemedisplay}

\begin{verbatim}
;; The cut-and-paste approach above seems wrong,
;; especially since the functions are so similar.
;; What if we wrote one function and made the operator
;; an argument?
\end{verbatim}

\newpage

\begin{schemedisplay}
(define (nums-filter pivot nl COMPARATOR)
  (cond
    [(empty? nl) empty]
    [(cons? nl)
     (cond
       [(COMPARATOR (first nl) pivot)
        (cons (first nl) (nums-filter pivot (rest nl) COMPARATOR))]
       [else
        (nums-filter pivot (rest nl) COMPARATOR)])]))
\end{schemedisplay}

\begin{verbatim}
; This works! Here's the contract for num-filter:
;   nums-filter : num num-list (num num -> boolean) -> num-list
\end{verbatim}


\begin{schemedisplay}
(define (nums-le pivot nl)
  (nums-filter pivot nl <=))

(define (nums-ge pivot nl)
  (nums-filter pivot nl -ge))
\end{schemedisplay}

\begin{schemedisplay}
$>$(nums-le 5 empty) 
empty
$>$(nums-le 5 (cons 1 empty))
(cons 1 empty)
$>$(nums-le 5 (cons 5 empty))
(cons 5 empty)
$>$(nums-le 5 (cons 6 empty))
empty
$>$(nums-le 5 (cons 1 (cons 6 (cons 5 empty))))
(cons 1 (cons 5 empty))
\end{schemedisplay}

\begin{verbatim}
; Now the main function.
;  quicksort : num-list -> num-list
;  returns a list containing the same numbers, sorted in ascending order
\end{verbatim}

\begin{verbatim}
; Template
;(define (quicksort nl)
;  (cond
;    [(empty? nl) ...]
;    [(cons? nl) ... (first nl) ...
;                ... (quicksort (rest nl)) ...]))

; This is a dead end:
; (quicksort (rest nl)) must return a sorted list,
; so the only choice is to insert (first nl) --
; which leads to some variant of insertion sort
\end{verbatim}

\newpage

\begin{schemedisplay}
;; Another idea
(define (quicksort nl)
  (cond
    [(empty? nl) ...]
    [else
     (combine-parts
      ... (quicksort (nums-le nl)) ...
      ...
      ... (quicksort (nums-ge nl)) ...)]))
\end{schemedisplay}

\begin{schemedisplay}
;; $>>$ BUGGY $<<$ version:
(define (quick-sort nl)
  (cond
    [(empty? nl) empty]
    [else
     (combine-parts
      (quick-sort (nums-le (first nl) nl))
      (first nl)
      (quick-sort (nums-ge (first nl) nl)))]))
;; Doesn't terminate!
\end{schemedisplay}

\begin{schemedisplay}
;; correct version (if input contains distinct numbers):
(define (quick-sort nl)
  (cond
    [(empty? nl) empty]
    [else
     (combine-parts
      (quick-sort (nums-le (first nl) (rest nl)))
      (first nl)
      (quick-sort (nums-ge (first nl) (rest nl))))]))
\end{schemedisplay}

\begin{schemedisplay}
(define (combine-parts le-nums pivot ge-nums)
  (append le-nums (cons pivot ge-nums)))
\end{schemedisplay}

\begin{schemedisplay}
$>$ (quick-sort empty) 
empty
$>$ (quick-sort (cons 1 empty))
(cons 1 empty)
$>$(quick-sort (cons 1 (cons 2 empty)))
(cons 1 (cons 2 empty))
$>$(quick-sort (cons 2 (cons 1 empty)))
(cons 1 (cons 2 empty))
$>$(quick-sort (cons 2 (cons 1 (cons 3 empty))))
(cons 1 (cons 2 (cons 3 empty)))
$>$(quick-sort (cons 11 (cons 8 (cons 14 (cons 7 empty)))))
(cons 7 (cons 8 (cons 11 (cons 14 empty))))
\end{schemedisplay}

\newpage

\begin{verbatim}
;; Recipe for generative recursion: 
;; The template is of the form
(define (func problem)
  (cond
   [(trivially-solveable? problem)  ...]
   [else
    (combine-solutions
     (func (generate-problem ...))
     ...
     (func (generate-problem ...)))]))

;; Another exaple of generative recursion: binary sort
\end{verbatim}

\begin{verbatim}
;; find-root num num (num -> num) -> num
;; Purpose: finds number n between a and b such that
;;  (f n) is within TOLERANCE of zero. f is an operator,
;;  (f a) must be negative, and (f b) must be positive.
\end{verbatim}
\begin{schemedisplay}
(define TOLERANCE 0.00005)
(define (find-root a b f)
  (cond
    [(<= (abs (- b a)) TOLERANCE) ...]
    [else
     (combine-parts
      (find-root a (/ (+ a b) 2) f)
      (find-root (/ (+ a b) 2) b f))]))
\end{schemedisplay}

\begin{schemedisplay}
(define (find-root a b f)
  (cond
    [(<= (abs (- b a)) TOLERANCE) a]
    [else
     (cond
      [(positive? (f (/ (+ a b) 2)))
       (find-root a (/ (+ a b) 2) f)]
      [else
       (find-root (/ (+ a b) 2) b f)])]))

$>$ (find-root -0.5 0.5 sin)
0
$>$ (find-root -2 0 cos)
-1.5712890625
\end{schemedisplay}

\newpage

\begin{verbatim}
;; We can make the repeated midpoint calculation,
;;  (/ (+ a b) 2), easier to read by using a local
;;  definition:
\end{verbatim}

\begin{schemedisplay}
(define (find-root a b f)
  (cond
    [(<= (abs (- b a)) TOLERANCE) a]
    [else
     (local [(define mid (/ (+ a b) 2))]
       (cond
         [(positive? (f mid))
          (find-root a mid f)]
         [else
          (find-root mid b f)]))]))
\end{schemedisplay}

\begin{verbatim}
;; What about roots for functions we define?
;; Let's try a polynomial: 4n^2 - 2n - 10
\end{verbatim}

\begin{schemedisplay}
(define (poly n)
  (+ (* 4 n n)
     (* -2 n)
     -10))

$>$ (poly 1)
-8
$>$ (poly 10)
370
$>$ (find-root 1 10 poly)
1.8507537841796875
\end{schemedisplay}

\begin{verbatim}
;; In other words: FUNCTIONS ARE VALUES
\end{verbatim}

\newpage

\begin{verbatim}
;; A more general number-filtering program:
;;   nums-filter2 : (num -> boolean) num-list -> num-list
\end{verbatim}

\begin{schemedisplay}
(define (nums-filter2 CHECKER nl)
  (cond
    [(empty? nl) empty]
    [(cons? nl)
     (cond
       [(CHECKER (first nl))
        (cons (first nl) (nums-filter2 CHECKER (rest nl)))]
       [else
        (nums-filter2 CHECKER (rest nl))])]))

(define (nums-le pivot nl)
  (local [(define (le-pivot n)
            (<= n pivot))]
    (nums-filter2 le-pivot nl)))
\end{schemedisplay}

\begin{verbatim}
;; Step through the first few evaluations above (using the
;;  new reduction rule for `local' expressions)

;; How about filtering animals?

;; animal-filter : (animal -> bool) animal-list -> animal-list
\end{verbatim}

\begin{schemedisplay}
(define (animal-filter CHECKER al)
  (cond
    [(empty? al) empty]
    [(cons? al)
     (cond
       [(CHECKER (first al))
        (cons (first al) (animal-filter CHECKER (rest al)))]
       [else
        (animal-filter CHECKER (rest al))])]))

(define (boa-filter al)
  (animal-filter boa? al))
\end{schemedisplay}

\begin{verbatim}
;; See any difference in the above filter programs, beside the names?
\end{verbatim}

\newpage

\begin{verbatim}
;; Let's just call it "filter"

;; filter : (T -> bool) T-list -> T-list
\end{verbatim}

\begin{schemedisplay}
(define (filter f? l)
  (cond
    [(empty? l) empty]
    [(cons? l)
     (cond
       [(f? (first l))
        (cons (first l) (filter f? (rest l)))]
       [else
        (filter f? (rest l))])]))
\end{schemedisplay}

\begin{verbatim}
;; This function is so useful that it's built into Scheme.

;; Another common thing we do to lists is convert each element
;;   (e.g., convert DM to dollars)

;; map : (T1 -> T2) T1-list -> T2-list
\end{verbatim}

\begin{schemedisplay}
(define (map f l)
  (cond
    [(empty? l) empty]
    [(cons? l) (cons (f (first l))
                     (map f (rest l)))]))

(define (convert rate l)
  (local [(define (convert-one n)
            (* rate n))]
     (map convert-one l)))

$>$ (convert 1.25 (list 4 10))
(list 5 12.5)
\end{schemedisplay}

\begin{verbatim}
;; Functions as values give us even more power: we can
;;  write functions that return functions!
;; This is useful for example, in using Newton's method
;;  to calculate roots. Given a function, we need its
;;  derivative...
\end{verbatim}

\newpage

{\tt ; square : num -$>$ num }
\begin{schemedisplay}
(define (square x)
  (* x x))
\end{schemedisplay}

{\tt ; d/dx-square : num -$>$ num}
\begin{schemedisplay}
(define EPSILON 0.01)
(define (d/dx-square n)
  (/ (- (square (+ n EPSILON)) (square (- n EPSILON)))
     (* 2 EPSILON)))
\end{schemedisplay}

{\tt ; d/dx-at-n : (num -$>$ num) num -$>$ num }
\begin{schemedisplay}
(define (d/dx-at-n square n)
  (/ (- (square (+ n EPSILON)) (square (- n EPSILON)))
     (* 2 EPSILON)))
\end{schemedisplay}

{\tt ; d/dx-at-n : (num -$>$ num) num -$>$ num }
\begin{schemedisplay}
(define (d/dx-at-n f n)
  (/ (- (f (+ n EPSILON)) (f (- n EPSILON)))
     (* 2 EPSILON)))
\end{schemedisplay}

{\tt ; d/dx-square : num -$>$ num }
\begin{schemedisplay}
(define (d/dx-square n) (d/dx-at-n square n))
\end{schemedisplay}

\begin{verbatim}
; But we really want d/dx : (num -> num) -> (num -> num)
\end{verbatim}
\begin{schemedisplay}
(define (d/dx f)
  (local [(define (d/dx-f n)
            (/ (- (f (+ n EPSILON)) (f (- n EPSILON)))
               (* 2 EPSILON)))]
      d/dx-f))
\end{schemedisplay}

\begin{verbatim}
; Why can't we return a function wihout using `local'?
; After all, we can write down numbers, positions, lists, and armadillos directly.
; Functions? We would like to be able to write something like he following.
\end{verbatim}

\begin{schemedisplay}
(define (d/dx f)
  (function (n) ;; $<$- a function of one argument, n
     (/ (- (f (+ n EPSILON)) (f (- n EPSILON))) ;; $<$- body
        (* 2 EPSILON))))
\end{schemedisplay}

\newpage

\begin{verbatim}
; Almost --- just use `lambda' instead of `function':
\end{verbatim}

\begin{schemedisplay}
(define (d/dx f)
  (lambda (n)
     (/ (- (f (+ n EPSILON)) (f (- n EPSILON)))
        (* 2 EPSILON))))
\end{schemedisplay}

\begin{verbatim}
; Why `lambda'? Tradition.

; We can use anonymous functions when calling d/dx, too:
\end{verbatim}

\begin{schemedisplay}
(define d/dx-square (d/dx (lambda (x) (* x x))))
(define d/dx-cube (d/dx (lambda (x) (* x x x))))
\end{schemedisplay}

\begin{verbatim}
; In fact, the `define' form you've learned for functions is just a
;  shorthand for `lambda'-valued abbreviations:
\end{verbatim}

\begin{schemedisplay}
(define square (lambda (x) (* x x)))
;; same as
(define (square x) (* x x))
\end{schemedisplay}

\begin{verbatim}
; convertDM is now a one-liner:
\end{verbatim}
\begin{schemedisplay}
(define (convertDM l) (map (lambda (x) (* x 1.82)) l))
\end{schemedisplay}

\begin{verbatim}
; Of course, abbreviations are just a convenience. It turns out that
; we could simplify the language by getting rid of `define', and
; just having `lambda' plus application forms. We don't even need numbers
; to do useful things, because we could just use `(lambda (x) (x))' 
; as a long name for `1', etc.

; Our claim has been that computing has nothing to do with computers.
; Now, by making our language more powerful, then throwing out
; redundant syntax, we have arrived at a brand new language --- 
; at least, it was brand new in the 1930's! Alonzo Church devised this
; Lambda Calculus long before mechanical computers were widely available.
\end{verbatim}

\end{document}
